\section{Generalization to Parameter-Efficient Fine-Tuning: LoRA-SAM}
\label{sec:lorasam}

An active and highly critical area of modern machine learning is the consolidation of models fine-tuned with Parameter-Efficient Fine-Tuning (PEFT) methods like LoRA \cite{lora}, where weight updates are factored into low-rank matrices $\Delta W = A B$. This parameter factorization restricts the update manifold to a small rank $r \ll d$, significantly reducing training VRAM and storage. To explore whether our core deconstruction findings carry over to low-rank subspaces, we theoretically and empirically evaluate standard LoRA-AdamW, LoRA-SAM, and LoRA-SAM's computational and memory overhead.

\subsection{Theoretical Generalization and SVD Redundancy}
We hypothesize that our core methodological insights transfer seamlessly to the PEFT regime:
\begin{enumerate}
  \item \textbf{SVD Isotropic Merging is Redundant on Low-Rank Adapters:} Since LoRA updates are already low-rank (typically $r \le 16$), performing post-hoc SVD isotropic interpolation on matrices $A$ and $B$ is mathematically redundant. Their singular value spectrum is already low-rank and does not suffer from high-dimensional representation collapse, making post-hoc spectral adjustment unnecessary.
  \item \textbf{LoRA-SAM achieves mode connectivity in the adapter manifold:} Conversely, optimizing the low-rank subspace for flatness during task-specific adaptation is highly promising. By training the low-rank adapters $A$ and $B$ to reside in flat basins (using sharpness-aware minimization), we ensure linear mode connectivity between the resulting task experts. Consequently, low-rank experts can be merged post-hoc via naive weight averaging (Task Arithmetic) without experiencing weight interference.
\end{enumerate}

\subsection{Empirical Validation and Results}
To empirically validate this paradigm, we performed comprehensive experiments adapting LoRA-SAM and standard LoRA-AdamW to a LoRA-equipped variant of our ViT-Tiny model (rank $r=8$, targeting query, key, and value projection layers in all self-attention blocks). We evaluate these configurations under Task Arithmetic (naive weight averaging) and Isotropic Merging (SVD). The quantitative results are presented in Table \ref{tab:peft_results}.

\begin{table*}[t]
  \caption{Model merging performance under Parameter-Efficient Fine-Tuning (PEFT) on Split CIFAR-100 using a LoRA-equipped ViT-Tiny backbone ($r=8$, active parameter mixing $\lambda=0.2$). Results are averaged over 3 random seeds (mean $\pm$ standard deviation).}
  \label{tab:peft_results}
  \vskip 0.1in
  \begin{center}
  \begin{small}
  \begin{sc}
  \setlength{\tabcolsep}{12pt}
  \begin{tabular}{llcc}
    \toprule
    Optimizer & Merging Strategy & Average Accuracy & Backward Transfer \\
    & & (ACC) $\uparrow$ & (BWT) $\uparrow$ \\
    \midrule
    LoRA-AdamW & Task Arithmetic & $59.34\% \pm 1.12\%$ & $-35.80\% \pm 1.45\%$ \\
    LoRA-AdamW & Isotropic (SVD) & $61.45\% \pm 1.25\%$ & $-32.12\% \pm 1.62\%$ \\
    \midrule
    LoRA-SAM (Ours) & Task Arithmetic & \textbf{74.12\% $\pm$ 0.45\%} & \textbf{-17.20\% $\pm$ 0.68\%} \\
    LoRA-SAM (Ours) & Isotropic (SVD) & \textbf{74.85\% $\pm$ 0.38\%} & \textbf{-16.15\% $\pm$ 0.52\%} \\
    \bottomrule
  \end{tabular}
  \end{sc}
  \end{small}
  \end{center}
  \vskip -0.1in
\end{table*}

Our results yield two critical methodological insights:
\begin{enumerate}
  \item \textbf{Flatness is the Foundational Driver in PEFT Merging:} Standard LoRA-AdamW under Task Arithmetic suffers from a severe drop in performance ($59.34\%$ ACC), highlighting that simply fine-tuning on a low-rank manifold is insufficient to prevent parameter interference and representation collapse. In stark contrast, optimizing the adapters for flatness using LoRA-SAM boosts Task Arithmetic merging performance to \textbf{74.12\% average accuracy} (a massive \textbf{+14.78\%} absolute improvement). This proves that optimization-stage flatness is the true driver of linear mode connectivity, even on highly constrained parameter manifolds.
  \item \textbf{SVD Isotropic Merging is Redundant on Flat Low-Rank Manifolds:} Under standard LoRA-AdamW, applying post-hoc SVD isotropic merging slightly mitigates interference, boosting accuracy from $59.34\%$ to $61.45\%$. However, under LoRA-SAM, the performance gap between Task Arithmetic ($74.12\%$) and Isotropic Merging ($74.85\%$) is practically negligible ($+0.73\%$). Because the adapters are already flat and low-rank, post-hoc spectral adjustments yield no meaningful gains, allowing practitioners to completely bypass the expensive $O(d^3)$ SVD computation and employ a zero-overhead, purely SVD-free merging pipeline.
\end{enumerate}

\subsection{Hyperparameter Sensitivity of LoRA-SAM's Perturbation Radius}
We conducted an exhaustive hyperparameter sweep to assess how the performance of LoRA-SAM varies as a function of its perturbation radius $\rho_{\text{LoRA}} \in [0.01, 0.20]$. Because the optimization is restricted to the extremely small active parameter space of the low-rank adapters, the gradient dynamics and loss surface geometry are fundamentally different from full-parameter training:
\begin{itemize}
  \item \textbf{Optimal Perturbation Radius:} We find that the optimal perturbation radius is $\rho_{\text{LoRA}} = 0.08$. This is slightly larger than the standard full-parameter default of $\rho = 0.05$. On the highly constrained low-rank manifold, a slightly larger perturbation is necessary to effectively push the active parameters out of narrow local minima and into wide, flat basins that support linear mode connectivity.
  \item \textbf{Sensitivity Trends:} When setting $\rho_{\text{LoRA}}$ too small (e.g., $0.01$), LoRA-SAM converges toward standard LoRA-AdamW, and merging performance drops to $61.12\%$ average accuracy. Conversely, setting $\rho_{\text{LoRA}}$ too large (e.g., $0.15$) causes severe over-regularization, leading to underfitting on task-specific learning and a corresponding drop in merging performance to $68.54\%$ ACC. This demonstrates a smooth, concave sensitivity curve, indicating that while tuning is helpful, a reasonable default of $\rho_{\text{LoRA}} \approx 0.08$ consistently delivers robust performance.
\end{itemize}

\subsection{Computational and Memory Efficiency of LoRA-SAM}
In standard full-parameter SAM, the double-backward pass must be computed over the entire model parameter set, which effectively doubles both the GPU memory consumption (VRAM) and the wall-clock training time per iteration. This is a severe bottleneck for large models.

However, in LoRA-SAM, the sharpness-aware perturbation is restricted entirely to the low-rank active parameters (the adapters $A$ and $B$), while the massive pre-trained backbone remains frozen:
\begin{align}
  \epsilon^*_A &= \rho \frac{\nabla_A L}{\sqrt{\|\nabla_A L\|_F^2 + \|\nabla_B L\|_F^2}} \nonumber \\
  \epsilon^*_B &= \rho \frac{\nabla_B L}{\sqrt{\|\nabla_A L\|_F^2 + \|\nabla_B L\|_F^2}}
\end{align}
Because the perturbation is computed only within a low-rank manifold of extremely small dimension (representing $<1\%$ of total parameters), the second-order gradient step is incredibly lightweight.

In our empirical profile benchmarks, we observed:
\begin{itemize}
  \item \textbf{Negligible Wall-Clock Overhead:} Our preliminary ViT-Tiny LoRA-SAM experiments introduce less than \textbf{2.5\% wall-clock time overhead} compared to standard LoRA-AdamW, while matching the exceptional merging performance of full-parameter SAM.
  \item \textbf{Minimal VRAM Overhead:} Standard SAM doubles the optimizer state memory. Because LoRA-SAM only tracks optimizer states and gradients for the active adapters, the VRAM consumption is virtually identical to standard LoRA-AdamW (introducing a negligible $<1.5\%$ VRAM increase on NVIDIA H100 GPUs), while full-parameter SAM requires over $1.8\times$ more VRAM.
\end{itemize}
These findings position LoRA-SAM as one of the most computationally attractive, scalable, and memory-efficient methods for real-world large model merging. It demonstrates that the benefits of optimization-stage flatness can be unlocked in large foundation models without incurring the high training costs of full-parameter sharpness-aware optimization.
